%% fig-layers.tex -- Figure 1: the six-layer stack (full width).
%%
%% Read bottom-up: infrastructure carries kernel services; kernel services obey
%% a governance substrate; the substrate governs asset families (the nouns);
%% lifecycle modes (the verbs) act on those nouns; surfaces expose the result.
%% The left rail holds platform services that cut across all six layers rather
%% than occupying one; its dashed outline marks that permeability.
%%
%% Type scale: every label is \figH (8pt) or \figB (7pt) -- see preamble.tex. The
%% content was trimmed to fit that floor rather than the floor lowered to fit the
%% content, because everything that came out is restated in the prose.

\begin{figure}[tbp]
\centering
\ifcbexcalidraw
  \begin{cbwide}\cbdiagram{fig-layers}\end{cbwide}
\else
\begin{cbwide}
\begin{tikzpicture}[x=1cm,y=1cm]

  % ---- geometry ------------------------------------------------------------
  % All box widths derive from one content width, so a band stays flush when a
  % label changes length. Total width stays inside the 7.0in text block.
  \def\CX{2.62}          % content left edge
  \def\CW{12.34}         % content width
  \def\G{0.18}           % gutter between boxes
  \def\LX{15.12}         % layer-label column left edge

  % ---- L6 surfaces (4) -----------------------------------------------------
  \foreach [count=\i from 0] \t in {%
      {\textbf{React SPA}\\{\figB same origin}},
      {\textbf{HTTP API}\\{\figB \statRouteDecls{} routes}},
      {\textbf{\aria copilot}\\{\figB tool loop}},
      {\textbf{Headless}\\{\figB SSE, webhooks}}}
  {
    \pgfmathsetmacro{\w}{(\CW-3*\G)/4}
    \node[surface, anchor=north west, text width=\w cm-8pt, minimum height=0.96cm]
      at ({\CX+\i*(\w+\G)}, 0) {\t};
  }

  % ---- L5 lifecycle modes (6) ---------------------------------------------
  \foreach [count=\i from 0] \t in {%
      {\mode{Author}\\{\figB draft, validate}},
      {\mode{Test}\\{\figB bounded runs}},
      {\mode{Evaluate}\\{\figB score, judge}},
      {\mode{Calibrate}\\{\figB optimize}},
      {\mode{Release}\\{\figB gate, promote}},
      {\mode{Observe}\\{\figB trace, meter}}}
  {
    \pgfmathsetmacro{\w}{(\CW-5*\G)/6}
    \node[control, anchor=north west, text width=\w cm-8pt, minimum height=0.96cm,
          inner xsep=2pt] at ({\CX+\i*(\w+\G)}, -1.34) {\t};
  }

  % ---- L4 asset families (9) ----------------------------------------------
  \foreach [count=\i from 0] \t in {%
      {Prompt}, {Tool}, {Skill}, {MCP}, {Work-\\flow},
      {Know.\\base}, {Test\\set}, {Judge}, {Agent}}
  {
    \pgfmathsetmacro{\w}{(\CW-8*\G)/9}
    \node[asset, anchor=north west, text width=\w cm-8pt, minimum height=0.84cm,
          font=\figB, inner xsep=1.5pt] at ({\CX+\i*(\w+\G)}, -2.68) {\t};
  }

  % ---- L3 governance (5) ---------------------------------------------------
  \foreach [count=\i from 0] \t in {%
      {\textbf{Identity}\\{\figB \statScopes{} RBAC scopes}},
      {\textbf{Exec.\ policy}\\{\figB governed egress}},
      {\textbf{Evidence}\\{\figB gate verdicts}},
      {\textbf{Release}\\{\figB promote, revert}},
      {\textbf{Ledgers}\\{\figB audit, effects}}}
  {
    \pgfmathsetmacro{\w}{(\CW-4*\G)/5}
    \node[govern, anchor=north west, text width=\w cm-8pt, minimum height=0.96cm,
          inner xsep=2pt] at ({\CX+\i*(\w+\G)}, -3.90) {\t};
  }

  % ---- L2 kernel (7) -------------------------------------------------------
  \foreach [count=\i from 0] \t in {%
      {Config \&\\secrets}, {Persis-\\tence}, {Storage}, {Tool\\exec.},
      {Event\\bus}, {Observ-\\ability}, {Provider\\adapters}}
  {
    \pgfmathsetmacro{\w}{(\CW-6*\G)/7}
    \node[control, anchor=north west, text width=\w cm-8pt, minimum height=0.96cm,
          font=\figB, inner xsep=1.5pt] at ({\CX+\i*(\w+\G)}, -5.24) {\t};
  }

  % ---- L1 infrastructure (5) ----------------------------------------------
  \foreach [count=\i from 0] \s/\t in {%
      store/{Starlette\\{\figB ASGI host}},
      store/{PostgreSQL 17\\{\figB pgvector, AGE}},
      store/{Object store\\{\figB S3, local}},
      extern/{\mlflow{} 3.14+\\{\figB registry, traces}},
      async/{Transport\\{\figB NATS, Redis}}}
  {
    \pgfmathsetmacro{\w}{(\CW-4*\G)/5}
    \node[\s, anchor=north west, text width=\w cm-8pt, minimum height=0.96cm,
          inner xsep=2pt] at ({\CX+\i*(\w+\G)}, -6.58) {\t};
  }

  % ---- the "rests on" spine ------------------------------------------------
  \foreach \ya/\yb in {-1.00/-1.30, -2.34/-2.64, -3.56/-3.86,
                       -4.90/-5.20, -6.24/-6.54}
    { \draw[flow, draw=cbMuted] (\CX+\CW/2, \ya) -- (\CX+\CW/2, \yb); }

  % ---- layer labels (right column) ----------------------------------------
  \foreach \y/\num/\name/\gloss in {
      0.00/6/SURFACES/the interface,
     -1.34/5/{LIFECYCLE MODES}/the verbs,
     -2.68/4/{ASSET FAMILIES}/the nouns,
     -3.90/3/GOVERNANCE/the rules,
     -5.24/2/KERNEL/services,
     -6.58/1/INFRASTRUCTURE/the base}
  {
    \node[anchor=north west, align=left, text width=2.66cm, font=\figB,
          text=cbInk, inner sep=0pt] at (\LX, \y-0.08)
         {\hyphenpenalty=10000\exhyphenpenalty=10000
          \textbf{\num}\dt\textbf{\name}\\[-0.1ex]
          {\itshape\color{cbMuted}\gloss}};
  }

  % ---- cross-cutting platform services (left rail) -------------------------
  % Dashed outline: these services are not a layer -- every layer reaches them.
  \node[muted, densely dashed, anchor=north west, text width=1.98cm, align=left,
        font=\figB, minimum height=7.54cm, inner xsep=3pt] at (0.0, 0.0)
    {\textbf{PLATFORM}\\\textbf{SERVICES}\\[0.4ex]
     {\itshape not a layer ---}\\{\itshape every layer}\\{\itshape reaches them}\\[1.0ex]
     Evidence base\\[0.45ex]
     Evaluation\\[0.45ex]
     Calibration\\[0.45ex]
     Capability\\registry\\[0.45ex]
     Integration hub\\[0.45ex]
     Project scoping};

\end{tikzpicture}
\end{cbwide}
\fi
\caption{The \sys stack. Six layers, read bottom-up: \Linfra{} carries
\Lkernel{}; kernel services obey \Lgov{}; the substrate governs \Lfam{} --- the
nouns; \Lmode{} --- the verbs --- act on those nouns; and \Lsurf{} exposes the
result. The dashed left rail holds platform services that cut across all six
layers rather than occupying one. Adjacency in this diagram is \emph{capability
availability, not inheritance}: a family in \Lnum{4} obtains \Lnum{5} behaviour
and \Lnum{3} enforcement only by explicitly wiring them, which is why the
guarantees in \tabref{tab:families} differ per row instead of being uniform. That
distinction is the paper's central architectural claim (\secref{sec:arch}).}
\label{fig:layers}
\figkey
\end{figure}
